The perception implementation is built around the idea that high-quality
airborne vision is constrained as much by sensing geometry and data
movement as by the neural model itself. For this reason, the codebase
extends the original single-board camera model into a dual-camera system
with explicit support for runtime switching, per-camera geometry, and
high-resolution acquisition.

At the capture level, the cameras are integrated through the processor
camera subsystem and the MIPI CSI-2 path. This is an important
architectural choice because it prevents the image path from becoming
dominated by the kind of software-managed peripheral transfer overhead
that would be unacceptable for sustained real-time inference. The camera
stack is designed so that CSI and MIPI remain active even during
switching, which reduces disruption and helps preserve throughput. In
practical terms, the runtime exposes \texttt{sentai.camera.switch(id)}
while internally snapshotting the frame sequence and waiting until
enough fresh frames have arrived from the newly selected sensor to
guarantee a clean image. This is a small but important example of the
thesis that real embedded autonomy depends on disciplined treatment of
edge cases that benchmark-oriented prototypes often ignore.

The OV5640 integration itself is more substantial than a generic sensor
bring-up. Register control is handled over SCCB through the board I2C
interfaces, including model-ID verification, reset and power-down
sequencing, per-camera mirror and rotation adjustments, and
sensor-specific initialization on each side of the camera multiplexer.
On the video side, the datasheet-level interface still matters to the
architecture: the sensor family supports compact formats such as 8-bit
RGB565, but the host capture path on the RT1176 is engineered around
aligned 32-bit framebuffers and DMA-friendly pitches. As a result, the
implementation deliberately separates transport format concerns from
model-input format concerns instead of assuming that the detector can
consume the capture representation directly.

The implementation also recognizes that switching cameras is not useful
unless geometry travels with the switch. The perception pipeline
therefore maintains per-camera configuration for field of view, mount
pitch, mount roll, mount yaw, and ground-reference mode. This allows the
same runtime to support parallel cameras, perpendicular camera
placement, nadir sensing, or oblique views without hard-coding
assumptions about the sensor pose. In aerial robotics this is essential
because the operational meaning of a detection depends on how the camera
is mounted, not only on where the pixels appear.

The neural perception path itself is intentionally deployment-aware.
Frames are captured at the sensor side, then scaled and converted
through the hardware pixel pipeline before being prepared for the model
input tensor. This means the system can begin from richer camera
imagery, including workflows derived from native 1280-wide sensing,
while still feeding smaller quantized tensors to the detector. Such a
path is crucial for compact YOLO-family models, which often live in
tension between the need for more spatial detail and the strict memory
limits of MCU deployment.

One fine but important engineering point is that the PXP stage is not
used only for resizing. The camera receiver exposes raw host buffers
whose layout is convenient for aligned capture, cache maintenance, and
bus transfer, but not for direct detector input. The runtime therefore
uses PXP to transform XRGB-style capture surfaces into packed RGB output
with the exact detector dimensions, and then quantizes in place if the
TPU model expects int8 input. This avoids repeated software byte
shuffling, avoids ambiguities caused by aligned scanline storage, and
turns pixel-format conversion into a deterministic hardware stage of the
pipeline.

Above the detector, the implementation adds a custom tracking layer
referred to in code as SentAI-SORT. The design borrows the logic of
modern tracking-by-detection systems while adapting it to the realities
of a Cortex-M7-class host. The tracker combines an 8-state Kalman
formulation, a ByteTrack-style two-stage association process, and
compact HSV histogram features. Instead of relying on expensive
feature-based motion compensation, it uses IMU-informed compensation
derived from camera motion. This substitution is not only
computationally cheaper, but also better aligned with the sensing
resources already present on the platform.

The tracker is further extended with ground-plane projection and
optional GPS-referenced localization. Once altitude, heading, IMU
attitude, and camera geometry are known, detections can be mapped beyond
image coordinates into a spatial frame meaningful for the drone and for
downstream consumers. This is particularly important when the runtime is
used not only to detect objects, but also to generate position-aware
telemetry or to feed higher-level navigation policies.

The implementation therefore supports a wider claim than simple onboard
detection. It demonstrates that a small embedded AI platform can host a
complete airborne perception chain: dual high-resolution sensing,
viewpoint switching, quantized detection, lightweight tracking, motion
compensation, and spatial projection. In combination with model
switching, this also creates the conditions for multi-model operation in
which different detectors or visual policies can be applied according to
mission phase, camera viewpoint, or scene constraints.

Finally, the perception stack is complemented conceptually by the custom
YOLO model and the synthetic open dataset discussed elsewhere in the
paper. Although those assets are external to the code implementation,
they are consistent with the implementation logic described here: the
runtime is engineered so that compact detection models trained for nadir
and oblique drone views can be loaded dynamically and used as part of a
practical onboard autonomy workflow.
